\scanpage{218}

\margin{\unsure{$4$|$5$|$'$}}

% The long strokes cancel the draft down to the closing bracket. The last two lines lie below them.
\deleted{associe le foncteur $F$ sur $D_S$ défini, pour $\unsure{\tau|\mathscr T}=$}

\diagram{\(\left(\begin{tikzcd}[ampersand replacement=\&,column sep=small,row sep=tiny]
\deleted{S}\&\deleted{T_0}\arrow[l,"\deleted{\varphi_0}"']\arrow[d,"\deleted{i}"]\\
\&\deleted{T}
\end{tikzcd}\right)\)}

\deleted{par $F(\deleted{\unsure{\tau_0}}\tau)=$ ens des hom $u:T\unsure{\text{—}|\to}U$ tels que $ui$ se factorise en $ju_0$, \inserted{tel}\deleted{\unsure{et}} que $\psi_0u_0=\varphi_0$. (NB. non ensembl\unsure{.})}

\diagram{\begin{tikzcd}[ampersand replacement=\&,column sep=small,row sep=small]
\deleted{S}\&\deleted{U_0}\arrow[l,"\deleted{\psi_0}"']\arrow[d,"\deleted{j}"']\&\deleted{T_0}\arrow[ll,bend right=18,"\deleted{\varphi_0}"']\arrow[l,dashed,"\deleted{u_0}"']\arrow[d,"\deleted{i}"]\\
\&\deleted{U}\&\deleted{T}\arrow[l,dashed,"\deleted{u}"']
\end{tikzcd}}

\deleted{[Si les voisinages infinitésimaux de tous ordres \inserted{$U^{(n)}$} de $U_0$ ds $U$ existent (p.ex. $j$ une imm) et si $\psi_0$ se relève, pour tt $n$, en $U^{(n)}\to S$ un \unsure{$R$-morphisme} à $\Sigma$, alors le foncteur $F$ qu'on vient de décrire est ind-représentable par \inserted{la famille} des}

\diagram{\(\left(\begin{tikzcd}[ampersand replacement=\&,row sep=tiny,column sep=small]
\deleted{S}\&\deleted{U^{\deleted{\unsure{n}}}_0}\arrow[l,"\deleted{\psi_0}"']\arrow[d,"\deleted{j^{(n)}}"]\\
\&\deleted{U^{(n)}}
\end{tikzcd}\right)\)}
\deleted{,]}\deleted{\unsure{de}}

\deleted{On \unsure{prouve|présume} en tous cas, grâce à l'hypothèse sur $\unsure{\mathscr D|D}$, \deleted{faite au début de cette \unsure{remarque}} \deleted{\unsure{que}} que $F$ est bien un \uline{faisceau} sur $D_S$, i.e. un objet de $\widetilde D_{S}$. Bien entendu, ce faisceau n'est pas en général ds $(\widetilde D_S)_{\unsure{\mathrm{strp}|\mathrm{str}|\mathrm{rep}}}$ ! [Cependant, $=$ sans doute dans $(\widetilde C)_{\unsure{\mathrm{strp}}}$, et $=$ si $U=U_0=S$, il peut être un torseur sous un groupe de $(\widetilde C)_{\unsure{\mathrm{strp}}}$ !!].}

cas $U=U_0=S$, $\psi_0=j=\mathrm{id}$. On trouve le foncteur $\tau\longmapsto$ ens des $\varphi$ qui relèvent $\varphi_0$.

% archive sheet number: 217
